\documentclass[12pt]{article}
\usepackage[margin=1in]{geometry}

% IPES 2026 Abstract — 500 words max
% Topic label: Money & Finance (or Other)

\begin{document}

\textbf{Who Pays, Who Punishes? Pandemic Wealth Shocks and Electoral Accountability Across Democracies}

\bigskip

Did the economic costs of the COVID-19 pandemic shape electoral accountability, and did the distribution of those costs across the wealth spectrum determine which incumbents were punished? We address this question using harmonized household wealth and income microdata from the Luxembourg Wealth Study (LWS) and Luxembourg Income Study (LIS), covering eight democracies that held elections during or shortly after the pandemic: the United States (2020), Canada (2021), the United Kingdom (2021), Japan (2021), Norway (2021), South Korea (2022), Denmark (2022), and Spain (2023).

Our empirical strategy proceeds in two stages. First, using LWS balance sheet data with pre-pandemic (2017--2019) and post-pandemic (2021--2023) waves, we measure how the pandemic redistributed wealth within each country. We decompose changes in household net worth by asset class---housing, equities, deposits, debt---and income quintile, identifying which segments of each country's wealth distribution bore the largest losses. We apply network-based community detection to household asset composition vectors across all 24 LWS countries, replacing the conventional stock-market-based versus bank-based typology with a continuous measure of financial structural similarity, and construct a country-level \textit{base exposure index} measuring the degree to which pandemic wealth shocks were concentrated among the incumbent party's electoral base.

Second, we test whether cross-national variation in base exposure predicts the magnitude of anti-incumbent voting swings in pandemic-era elections. Using subnational election data at the county, constituency, or precinct level, we estimate the within-country relationship between local wealth vulnerability and vote change, then pool estimates across countries using network-derived financial structure groupings. We assess robustness through subnational spatial network models that account for geographic dependence in both wealth shocks and voting behavior.

Our identification strategy exploits the fact that pre-pandemic wealth composition---the share of household assets held in equities versus housing versus deposits---is structurally determined by long-run institutional factors that predate the pandemic. We construct a network-weighted shift-share instrument: pre-pandemic asset composition provides the shares, while the shift is mediated through cross-border financial linkage networks constructed from BIS and IMF bilateral banking and portfolio investment data. This isolates the component of wealth change attributable to structural financial exposure and cross-border contagion rather than domestic policy responses.

Preliminary evidence from the United States supports this framework. Using individual-level data from the Cooperative Election Study and county-level election returns, we find that voters who experienced personal economic costs---particularly job loss and bereavement---were significantly more likely to defect from the incumbent, and this effect is concentrated among the incumbent's prior supporters. A shift-share instrumental variables analysis confirms the causal interpretation. The present paper extends this single-country finding to a comparative framework that integrates network methods with causal inference.

This paper contributes to the IPE literature on the domestic political consequences of economic shocks by showing that the \textit{distributional} incidence of crisis costs---not merely their aggregate magnitude---determines electoral accountability. The findings illuminate how financial globalization, by embedding household portfolios in cross-border networks of asset exposure, conditions the political response to international crises.

\end{document}
